\documentclass[11pt,a4paper]{article}
\usepackage[top=2cm,bottom=2cm,left=2.5cm,right=2.5cm]{geometry}
\usepackage{parskip}
\usepackage{microtype}
\usepackage{eurosym}
\usepackage{enumerate}
\setlength{\parskip}{4pt}

\begin{document}
\pagestyle{empty}

{\large\textbf{RDC\_Peer\_Review 27.06.2026}}

Peer Review by Hari Krishna Mohan Raj

\bigskip

\textbf{\large Paper Title: Data-Driven Makeup Detection via TOF Cameras in Biometrics}

\textbf{\large Author: Aidin Masroor, Sebastian Houben, Ralph Breithaupt, Markus Rohde}

\bigskip

\textbf{Intended Conference: IEEE Transactions on Pattern Analysis and Machine Intelligence (TPAMI) 2027}

\bigskip

Rating: (Exceptional $|$ Major $|$ Good $|$ Minor $|$ Questionable $|$ None $|$ Previously Published)

\begin{enumerate}[(i)]

\item \textbf{Relevance of the use of AI in the application/context: Major}

The idea of using a deep learning model to figure out what a face should look like
from a depth image makes sense here, since you cannot really write a simple rule for
that. The authors tried three different model types and picked the one that worked
best, which shows they thought it through rather than just going with the first option.

\item \textbf{Significance of contribution: Good}

Using a ToF camera to catch people wearing makeup to fool face recognition is a
fresh idea and the paper says no one else has done it this way before. It also seems
cheaper than the setups used in other work, which is a practical plus. The results
look promising even if the study is still at an early stage.

\item \textbf{Technical quality of the intelligent solution: Good}

The model seems to work reasonably well based on the images shown and the similarity
score plots. The authors also tested a few different model architectures before
settling on U-Net, which adds some confidence to the choice. There is a missing
reference marked as \texttt{??} in the text which should be fixed before submitting.

\item \textbf{Clarity and quality: Good}

The paper is easy enough to follow and the figures help a lot in understanding what
is going on. Some parts of the writing are a bit rough and there are a few spelling
mistakes here and there. A quick proofread would clean things up and make it read
more polished overall.

\end{enumerate}

\medskip

\textit{Reviewer's confidence: 4 / 5}

\textbf{Overall Evaluation:} Overall, I would rate this as an accept. The idea is
interesting and the approach seems to work based on what is shown. It is a fairly
new way of tackling the problem and the paper is understandable. Some small things
like spelling mistakes and a missing reference should be fixed.

\end{document}
